\section{Estructura del documento}

El presente documento se encuentra dividido en seis capítulos, además de los anexos que complementan la información presentada.

El Capítulo 1, \textit{Contexto y objetivos}, plantea el problema de investigación, formula el objetivo general y los objetivos específicos del trabajo, y delimita el alcance técnico y evaluativo del prototipo a desarrollar.

El Capítulo 2, \textit{Marco Referencial}, presenta el glosario de términos empleados a lo largo del documento, el marco teórico relacionado con los fundamentos de la interpretación de lenguajes de programación, las teorías del aprendizaje aplicadas a la enseñanza de compiladores y la relación conceptual entre gramática, AST y ambiente de ejecución, así como los antecedentes identificados mediante una búsqueda sistemática de literatura.

El Capítulo 3, \textit{Diagnóstico de dificultades conceptuales}, describe la metodología empleada para identificar las dificultades de aprendizaje en la asignatura FLP, presenta el análisis de las encuestas aplicadas a estudiantes y docentes, y expone la selección de los conceptos base que orientan el diseño del prototipo a partir de la revisión del microcurrículo de la asignatura.

El Capítulo 4, \textit{Diseño del prototipo FLP Viewer}, especifica los requisitos funcionales y no funcionales del sistema y describe las decisiones de diseño adoptadas, incluyendo la arquitectura del sistema, la selección de tecnologías, los prototipos de interfaz de usuario y las estrategias de visualización.

El Capítulo 5, \textit{Implementación del prototipo FLP Viewer}, documenta la implementación de los componentes del sistema, el flujo de ejecución principal, el proceso de despliegue en el entorno de pruebas y la verificación de los requisitos funcionales definidos en el Capítulo 4.

El Capítulo 6, \textit{Conclusiones}, sintetiza los resultados alcanzados en contraste con los objetivos planteados y presenta las conclusiones generales del trabajo de grado.

Finalmente, los anexos recogen el microcurrículo de la asignatura FLP, los instrumentos y resultados de las encuestas aplicadas a estudiantes y docentes, el protocolo de búsqueda sistemática de literatura y las especificaciones de requisitos del sistema.